\section{Zadanie 3}
Niech $B\subseteq\Omega$ i niech $\sum_B = \{ A\subseteq\Omega :B\subseteq A$ lub $B\cap A = \emptyset\}$. Udowodnij, że $\sum_B$  jest $\sigma$-algebrą.
\subsection*{Rozwiązanie}
Przejdźmy kolejno po poszczególnych przypadkach.
\begin{enumerate}
\item Zbiór $\emptyset\in \sum_B$ ponieważ $B\cap\emptyset=\emptyset$. Zbiór $\Omega\in \sum_B$ ponieważ $B\subseteq\Omega$.
\item Weźmy $E\in\sum_B$ i zauważmy że jeśli $E$ spełnia warunek, że $E\cap B = \emptyset$ to $B\subseteq\Omega\setminus E$. Jeśli zaś zachodzi, że $B\subseteq E$ to zachodzi $(\Omega\setminus E)\cap B = \emptyset$.
\item Jeżeli $\exists_i$ takie że $E_i$ spełnia warunek $B\subseteq E_i$ to wtedy $B\subseteq \bigcup_{i=1}^n E_i$ zatem $\bigcup_{i=1}^n E_i \in \sum_B$. Jeżeli taki zaś nie istnieje to dla każdego $i$ zachodzi $E_i\cap B=\emptyset$ zatem $\bigcup_{i=1}^n E_i \cap B=\emptyset$ czyli $\bigcup_{i=1}^n E_i\in\sum_B$.
\item analogicznie jak 3 podpunkt bo nie korzystamy tam z tego że suma jest skończona
\end{enumerate}
